\scanpage{101}
% Two long oblique strokes run down across the typescript; the first begins at « pour ».
\deleted{pour lesquels l'expression $(*)$ est définie.}

\deleted{fermé $Z_s$ de $X_s$, où $Z$ est le sous-schéma fermé de $X$ adhérence \deleted{\unsure{de}}schématique de $Z_t$ (de sorte que $Z$ est le sous-schéma fermé intègre de $X$ défini par l'adhérence topologique de $Z_t$). \deleted{Je dis} \inserted{Il faudrait montrer} que cet homomorphisme passe au quotient pour définir un homomorphisme d'anneaux gradués}
\[
\deleted{\mathord{:}\ A^{\bullet}(X_t)\qquad A^{\bullet}(X_s)\quad .}\tag*{(7.11.2)}
\]

\deleted{Pour montrer que (7.10.1) transforme cycles linéairement équivalents à zéro en cycles linéairement équivalents à zéro, on introduit le schéma $C=P^1_S$, les deux sections $x_{\unsure{o|0}}$ et $x_1$ de $C$ sur $S$ définies par $0$ et $1$, et on se rappelle que le groupe des cycles \unsure{linéairement|linéairment} équivalents à zéro sur $X_t$ est engendré par les cycles}
\[
\deleted{Z_t(x_{1t})\ \deleted{\unsure{=|\times}}\unsure{-|=}\ Z_t(x_{ot})\ ,}\tag*{$(*)$}
\]
\deleted{où $Z_t$ parcourt le\deleted{s cycles} groupe des cycles sur $X_t\mathbin{\mathrm{x}}_{k(t)}C_t$\handwritten{\unsure{!}}. Or un calcul d'intersections facile (\hspace{3em}) montre que l'image par (7.11.1) du cycle $(*)$ n'est autre que}
\[
\deleted{Z_s(x_{1s})-Z_s(x_{os})\quad ,}\tag*{$(*\ *)$}
\]
\deleted{où $Z$ désigne encore l'adhérence schématique de $Z_t$ dans $X$, \uline{pourvu que l'expression} $(*\ *)$ \uline{soit définie}.}
% archive sheet number: 100
